Nuclear power plants continuously acquire neutron and instrument-channel measurements for surveillance and condition monitoring. These raw time series can reveal changes in detector behavior, signal integrity, and operating regimes, but exhaustively assigning reliable fault labels to long, multichannel archives requires context-sensitive expert review with operating and maintenance information. That effort is impractical at routine archive scale, and the CSV records used here consequently do not provide a uniform sample-level fault ground truth. Technical interpretation therefore continues to rely on authoritative online-monitoring and instrument-channel-assessment guidance~\citep{iaea2008olm1,iaea2008olm2,hashemian2011online}. Neutron-noise techniques provide a non-intrusive route to studying reactor and instrumentation phenomena~\citep{pazsit2010noise}, while projects such as CORTEX extend numerical simulation and validation capabilities for reactor monitoring~\citep{demaziere2022cortex}. Data-driven studies have demonstrated scenario-specific neutron-detector validation and nuclear-plant fault classification~\citep{tambouratzis2020neutron,elbordany2024fault}; however, the fixed signal constructions and predefined fault settings used in such studies are not always available in unlabeled plant archives.

Large language models (LLMs) can help organize heterogeneous observations and technical knowledge, but free-form generation is unsuitable as a decision authority in a safety-relevant setting. Retrieval-augmented generation (RAG)~\citep{lewis2020rag} conditions an answer on external evidence; in nuclear applications, its value depends on whether the retrieved manuals genuinely support the statement being made. Nuclear-domain studies use retrieval and factuality checks to improve the reliability of technical explanations~\citep{reeves2024reliability,lin2025cws}. Document-interpretation experiments identify failure modes including incorrect retrieval, unsupported interpretation, hallucination, and refusal~\citep{mapes2025uprate}, while RADIANT-LLM emphasizes provenance tracking, citation-backed responses, and human validation~\citep{ndum2026radiant}. A practical open question is how to connect measurement-side findings to that evidence layer without asking an LLM to infer a root cause directly from a raw CSV file.

This work addresses that question through an evidence-constrained diagnostic-assistance workflow. A tool-using agent observes unlabeled raw neutron-current series with minimal diagnostic priors and writes an auditable descriptive condition report. This report is a derived analysis artifact, not a ground-truth annotation: it records signal phenomena such as spikes, zero events, co-movement, lag, and data-health cues without promoting them to a fault class. It is then compressed into a structured alert summary that guides human-specified retrieval and citation-bearing explanation. The intended output is not an automatically confirmed fault diagnosis: it separates observed data phenomena, manual-supported candidate interpretations, and the engineering review that remains necessary to judge either.

The reliability of this workflow depends critically on retrieval. A manual library may need to include both a small, pre-designated authoritative set and a larger set of same-domain supplements. Adding relevant supplements can improve apparent nearest-neighbor similarity while reducing the priority of the authoritative sources specified by a deployment policy. This tension is especially important for bilingual technical terminology and for overlap-based chunking, where a useful passage may not have the identical chunk identifier used as a query probe.

The contributions are threefold. First, we develop a zero-prior, tool-using workflow with a persistent workspace contract for converting neutron-current CSV data into auditable descriptive observations, providing a scalable intermediate layer when exhaustive fault labeling is unavailable. Second, we introduce report-driven multi-stage RAG that maps structured observations to citable manual evidence. Third, we evaluate evidence retrieval under controlled core--supplement expansion, query language, chunking, and strict versus document-level relevance. The resulting evidence is used to characterize deployment boundaries rather than to claim autonomous fault identification. The complete transformation contract is specified in \cref{sec:system}.
